% =============================================================================
% ForensiQ — Gestão de Riscos com Controlos
% UC 21184 · Projeto de Engenharia Informática · UAb
% =============================================================================
\documentclass[a4paper,11pt]{article}

\usepackage[utf8]{inputenc}
\usepackage[T1]{fontenc}
\usepackage[top=2.5cm, bottom=2.5cm, left=3cm, right=2.5cm]{geometry}
\usepackage{booktabs}
\usepackage{longtable}
\usepackage{array}
\usepackage{parskip}
\usepackage{hyperref}
\usepackage{xcolor}
\usepackage{colortbl}
\usepackage{enumitem}

\hypersetup{
    colorlinks=true,
    linkcolor=black,
    urlcolor=blue,
    pdfauthor={João Rodrigues},
    pdftitle={ForensiQ — Gestão de Riscos com Controlos}
}

\newcolumntype{L}[1]{>{\raggedright\arraybackslash}p{#1}}

% =============================================================================
\begin{document}
% =============================================================================

\begin{center}
    {\LARGE\bfseries Gestão de Riscos com Controlos}\\[0.8em]
    {\large ForensiQ --- Plataforma Modular de Gestão de Prova Digital\\
    para First Responders}\\[1.2em]
    \begin{tabular}{ll}
        \textbf{Versão:} & 2.0 · 12 abril 2026 \\
        \textbf{Referência:} & ISO/IEC 27005:2018 (Risk management) \\
    \end{tabular}
\end{center}

\vspace{0.5em}
\hrule
\vspace{1.5em}

% =============================================================================
\section{Introdução}
% =============================================================================

Este documento estende a análise de riscos do projeto ForensiQ com uma matriz
formal de risco-para-controlo. Para cada risco identificado, descreve-se a
probabilidade, impacto, controlo(s) implementado(s) e estado atual de mitigação.

A metodologia segue os princípios de ISO/IEC 27005 para gestão de riscos de
segurança da informação, adaptados ao contexto de um projeto académico de
engenharia informática.

% =============================================================================
\section{Riscos do Projeto (Gestão)}
% =============================================================================

\subsection*{R01 — Tempo insuficiente para MVP completo}

\begin{itemize}[nosep]
    \item \textbf{Probabilidade:} Alta (> 75\%)
    \item \textbf{Impacto:} Alto (atraso crítico)
    \item \textbf{Controlo:} MVP-first rigoroso com prioridade 1-2-3 (MoSCoW).
          Funcionalidades avançadas apenas se núcleo sólido. Plano de entrega em fases.
    \item \textbf{Estado:} Mitigado
\end{itemize}

\subsection*{R02 — Dificuldade em aprender Django/DRF}

\begin{itemize}[nosep]
    \item \textbf{Probabilidade:} Média (25--75\%)
    \item \textbf{Impacto:} Alto
    \item \textbf{Controlo:} Código com revisão obrigatória. Começar pelo modelo
          de dados (ponto forte). Reutilização de padrões DRF. Documentação interna.
    \item \textbf{Estado:} Mitigado
\end{itemize}

\subsection*{R03 — Orientador rejeitar proposta}

\begin{itemize}[nosep]
    \item \textbf{Probabilidade:} Baixa (< 25\%)
    \item \textbf{Impacto:} Crítico (projeto não segue)
    \item \textbf{Controlo:} Enviar elementos base com margem antes do prazo.
          Incorporar feedback iterativamente. Reuniões semanais com orientador e tutor.
    \item \textbf{Estado:} Mitigado
\end{itemize}

\subsection*{R04 — Geolocalização não funcionar em demonstração}

\begin{itemize}[nosep]
    \item \textbf{Probabilidade:} Média
    \item \textbf{Impacto:} Baixo (fallback manual disponível)
    \item \textbf{Controlo:} Fallback: coordenadas manuais (texto livre).
          Testar com telemóvel real antes de demo. GPS simulado em dev.
    \item \textbf{Estado:} Aceite
\end{itemize}

\subsection*{R05 — Hash de integridade difícil de demonstrar}

\begin{itemize}[nosep]
    \item \textbf{Probabilidade:} Baixa
    \item \textbf{Impacto:} Baixo
    \item \textbf{Controlo:} Demonstração directa na BD: alterar campo invalida hash.
          Exportação PDF mostra hash. Testes automatizados para integridade.
    \item \textbf{Estado:} Mitigado
\end{itemize}

% =============================================================================
\section{Riscos de Segurança (Implementação)}
% =============================================================================

\subsection*{R06 — Alteração maliciosa de prova digital}

\begin{itemize}[nosep]
    \item \textbf{Probabilidade:} Baixa
    \item \textbf{Impacto:} Crítico (compromete integridade jurídica)
    \item \textbf{Controlo:} Imutabilidade em 3 camadas:
    \begin{enumerate}[nosep, leftmargin=*]
        \item Django: \texttt{read-only=True} no model
        \item API: validação no serializador DRF rejeita UPDATE
        \item PostgreSQL: \emph{triggers} rejeitam UPDATE a campos críticos
    \end{enumerate}
    Hash SHA-256 recalculado valida integridade em cada acesso.
    \item \textbf{Estado:} Mitigado
\end{itemize}

\subsection*{R07 — Acesso não autorizado a dados de outro agente (IDOR)}

\begin{itemize}[nosep]
    \item \textbf{Probabilidade:} Média
    \item \textbf{Impacto:} Alto (compromete confidencialidade)
    \item \textbf{Controlo:} Filtro automático de queryset por \texttt{agent-id}.
          Cada AGENT vê apenas suas ocorrências. Testes IDOR em endpoints críticos.
          Validação ao nível do serializador.
    \item \textbf{Estado:} Mitigado
\end{itemize}

\subsection*{R08 — Brute-force em credenciais JWT}

\begin{itemize}[nosep]
    \item \textbf{Probabilidade:} Média
    \item \textbf{Impacto:} Alto (acesso não autorizado)
    \item \textbf{Controlo:} Rate limiting (10 req/min anónimos). JWT com rotação
          automática e blacklist. Falha de autenticação registada em AuditLog.
          Senha com hash PBKDF2.
    \item \textbf{Estado:} Mitigado
\end{itemize}

\subsection*{R09 — Race condition na cadeia de custódia}

\begin{itemize}[nosep]
    \item \textbf{Probabilidade:} Baixa
    \item \textbf{Impacto:} Alto (inconsistência de estado)
    \item \textbf{Controlo:} \texttt{select-for-update()} na view para lock
          exclusivo. \texttt{transaction.atomic()} encapsula operação. Testes com
          concorrência.
    \item \textbf{Estado:} Mitigado
\end{itemize}

\subsection*{R10 — Perda de rastreabilidade de acessos}

\begin{itemize}[nosep]
    \item \textbf{Probabilidade:} Média
    \item \textbf{Impacto:} Alto (impossibilidade de auditoria)
    \item \textbf{Controlo:} Módulo AuditLog estruturado com: utilizador, ação,
          recurso, parâmetros, timestamp, IP, correlation ID, resultado. Auditoria
          consultável e descarregável. Impossível modificar histórico.
    \item \textbf{Estado:} Mitigado
\end{itemize}

% =============================================================================
\section{Riscos de Operação (Infraestrutura)}
% =============================================================================

\subsection*{R11 — Exposição de credenciais em logs ou URLs}

\begin{itemize}[nosep]
    \item \textbf{Probabilidade:} Média
    \item \textbf{Impacto:} Crítico (roubo de credenciais)
    \item \textbf{Controlo:} Sanitização de logs (password, token nunca em plaintext).
          JWT nunca em URL (header \texttt{Authorization} obrigatório). Admin Django
          em URL aleatória, não em \texttt{/admin}. Scan de código para secrets.
    \item \textbf{Estado:} Mitigado
\end{itemize}

\subsection*{R12 — Falha de disponibilidade durante defesa}

\begin{itemize}[nosep]
    \item \textbf{Probabilidade:} Média
    \item \textbf{Impacto:} Médio (demo não funciona)
    \item \textbf{Controlo:} Deploy em plataforma cloud com SLA. Monitoramento
          de uptime. Teste de carga antes de defesa. Documentação de rollback.
    \item \textbf{Estado:} Aceite
\end{itemize}

\subsection*{R13 — Injeção de SQL ou XSS em formulários}

\begin{itemize}[nosep]
    \item \textbf{Probabilidade:} Baixa
    \item \textbf{Impacto:} Crítico (compromete integridade de BD)
    \item \textbf{Controlo:} Django ORM previne SQL injection. Serialização DRF
          com validação de tipo. Template escapes HTML automaticamente. CSP header
          restritiva.
    \item \textbf{Estado:} Mitigado
\end{itemize}

% =============================================================================
\section{Tabela Resumida de Riscos}
% =============================================================================

\begin{longtable}{L{0.7cm} L{4cm} L{1.2cm} L{1.2cm} L{3.5cm} L{1.2cm}}
    \toprule
    \textbf{ID} & \textbf{Risco} & \textbf{Prob.} & \textbf{Imp.} & \textbf{Controlo Principal} & \textbf{Estado} \\
    \midrule
    \endhead
    R01 & Tempo insuficiente & Alta & Alto & MVP-first & Mitigado \\
    R02 & Curva Django & Média & Alto & Revisão código & Mitigado \\
    R03 & Rejeição proposta & Baixa & Crítico & Feedback iterativo & Mitigado \\
    R04 & GPS falha & Média & Baixo & Fallback manual & Aceite \\
    R05 & Demo de hash & Baixa & Baixo & Testes automatizados & Mitigado \\
    R06 & Alteração prova & Baixa & Crítico & Imutabilidade 3x & Mitigado \\
    R07 & IDOR & Média & Alto & Filtro queryset & Mitigado \\
    R08 & Brute-force & Média & Alto & Rate limiting & Mitigado \\
    R09 & Race condition & Baixa & Alto & select-for-update & Mitigado \\
    R10 & Perda auditoria & Média & Alto & AuditLog estruturado & Mitigado \\
    R11 & Exposição credenciais & Média & Crítico & Sanitização logs & Mitigado \\
    R12 & Downtime & Média & Médio & Cloud SLA & Aceite \\
    R13 & Injeção SQL/XSS & Baixa & Crítico & Django + CSP & Mitigado \\
    \bottomrule
\end{longtable}

% =============================================================================
\section{Escalas de Avaliação}
% =============================================================================

\textbf{Probabilidade:}
\begin{itemize}[nosep]
    \item Baixa: < 25\%
    \item Média: 25--75\%
    \item Alta: > 75\%
\end{itemize}

\textbf{Impacto:}
\begin{itemize}[nosep]
    \item Baixo: Atraso menor, sem impacto em entrega
    \item Médio: Atraso significativo ou funcionalidade comprometida
    \item Alto: Atraso crítico ou funcionalidade crítica em risco
    \item Crítico: Projeto não entregar ou compromisso legal/segurança
\end{itemize}

\textbf{Estado de Mitigação:}
\begin{itemize}[nosep]
    \item \textbf{Mitigado:} Controlo implementado, probabilidade residual < 10\%
    \item \textbf{Aceite:} Risco tolerável (baixa probabilidade ou baixo impacto)
    \item \textbf{Em monitorização:} Controlo parcial, monitorização contínua
\end{itemize}

% =============================================================================
\section{Plano de Monitorização}
% =============================================================================

\begin{longtable}{L{0.7cm} L{2.5cm} L{4cm} L{2.5cm}}
    \toprule
    \textbf{ID} & \textbf{Risco} & \textbf{Indicador} & \textbf{Frequência} \\
    \midrule
    \endhead
    R01 & Tempo & Progresso de requisitos (\%) & Semanal \\
    R07 & IDOR & Testes de IDOR, tentativas crossagent em AuditLog & Mensal \\
    R08 & Brute-force & Taxa de falhas de autenticação & Diário \\
    R10 & Auditoria & Integridade de AuditLog & Mensal \\
    R11 & Credenciais & Scan de código para secrets & Semanal \\
    R12 & Downtime & Uptime do servidor & Diário \\
    \bottomrule
\end{longtable}

% =============================================================================
\section{Histórico de Actualização}
% =============================================================================

\begin{longtable}{L{2cm} L{2.5cm} L{5.5cm} L{2cm}}
    \toprule
    \textbf{Data} & \textbf{Versão} & \textbf{Alteração} & \textbf{Estado} \\
    \midrule
    \endhead
    22 mar 2026 & 1.0 & Versão inicial com 5 riscos de projeto & --- \\
    12 abr 2026 & 2.0 & Adicionados 8 riscos de segurança/operação + controlos & Atualizado \\
    \bottomrule
\end{longtable}

\end{document}
